\documentclass{article}
\usepackage{graphicx} % Required for inserting images
%\usepackage[a4paper, total={6in, 8in}]{geometry}
\usepackage{amsmath, amssymb, empheq, mathtools, bbold, enumitem, hyperref, physics}
\usepackage{amsthm}
\usepackage[framemethod=TikZ]{mdframed}
%\usepackage{framed}
%\usepackage[amsthm,framed,thmmarks]{ntheorem}

\renewcommand{\bf}[1]{\textbf{#1}}
\renewcommand{\v}[1]{\boldsymbol{#1}}

%\newframedtheorem{frm-thm}{Theorem}
\newtheorem{theorem}{Theorem}
\newtheorem*{fact}{Fact}
\newtheorem*{note}{Note}
\newtheorem{definition}{Definition}

\surroundwithmdframed[]{theorem}
\newenvironment{frmthm}{
	\begin{mdframed}
		\begin{theorem}
		}{		
		\end{theorem}
	\end{mdframed}
}

\newcommand{\xlrightarrow}{\xrightarrow{\hspace{0.8cm}}}
\newcommand{\reals}{\mathbb{R}}

\newcommand{\matadd}{+_{\text{mat}}}
\newcommand{\matdot}{\cdot_{\text{mat}}}

\newcommand{\mxn}{m \times n}
\newcommand{\letmatrix}[2]{\mathcal{M}_{#1 \times #2}(\reals)}
\newcommand{\rowop}{\mathcal{R}}

\newcommand{\bmat}[1]{\begin{bmatrix}#1\end{bmatrix}}
\newcommand{\detmat}[1]{
	\left|\hspace{0.5em}
	\begin{matrix}
		#1
	\end{matrix}
	\hspace{0.5em}\right|
}
% Ch. 2: Linear Combinations and Linear Independence

\begin{document}
\setcounter{section}{-1}
\section{Preface}
\section[Vectors in R^n]{Vectors in $R^n$}
\begin{definition}[Vectors in $\reals^n$]
	Euclidean $n$-space, denoted by $\reals^n$, or simply $n$-space, is defined by
	\[ \reals^n = \left\{ \bmat{x_1\\x_2\\\cdots\\x_n} \vrule \; x_i \in \reals, \text{for} i = 1, 2, ..., n \right\} \]
	The entries of a vector are called the \textbf{components} of the vector.
\end{definition}
The vector $v = \left[\begin{smallmatrix} 1\\2 \end{smallmatrix}\right]$ is the directed line segment from the origin $(0, 0)$ to the point $(1, 2)$. The point $(0, 0)$ is the \textbf{initial point} and the point $(1,2)$ is the \textbf{terminal point}. The length of the vector is the length of the line segment from the initial point to the terminal point. For example, the length of $v = \bmat{1\\2}$ is $\sqrt{1^2 + 2^2} = \sqrt{5}$. A vector is unchanged if it is relocated elsewhere in the plane, provided that the length and direction remain unchanged. When the initial point of a vector is the origin, we say that the vector is in \textbf{standard position}.\par
Since vectors are matrices, two vectors are \textbf{equal} provided that their corresponding components are equal. The operations of addition and scalar multiplication are defined componentwise as they are for matrices.
\begin{definition}[Addition and Scalar Multiplication of Vectors]
	Let \bf{u} and \bf{v} be vectors in $\reals^n$ and $c$ a scalar.
	\begin{enumerate}[itemsep=0pt]
		\item The \bf{sum} of \bf{u} and \bf{v} is
		\[ \text{\bf{u} $+$ \bf{v}} = \bmat{u_1\\u_2\\ \cdots \\ u_n} + \bmat{v_1\\v_2\\ \cdots \\ v_n} = \bmat{u_1 + v_1\\u_2 + v_2 \\ \cdots \\ u_n + v_n} \]
		\item The \bf{scalar product} of $c$ and \bf{u} is
		\[ \text{$c$\bf{u}} = c\bmat{u_1\\u_2\\\cdots\\u_n} = \bmat{cu_1\\cu_2\\\cdots\\cu_n} \]
	\end{enumerate}
\end{definition}
\section{Linear Combinations}
\begin{definition}[Linear Combination]
	Let $S = \{ \vb{v}_1, \vb{v}_2, \dots, \vb{v}_k \}$ be a set of vectors in $\reals^n$, and let $c_1, c_2, \dots, c_k$ be scalars. An expression of the form
	\[ c_1\vb{v}_1 + c_2\vb{v}_2 + \cdots + c_k\vb{v}_k = \sum_{i=1}^{k}c_i\vb{v}_i \]
	is called a \bf{linear combination} of the vectors of $S$. Any vector \bf{v} that can be written in this form is also called a linear combination of the vectors of $S$.
\end{definition}


\end{document}